% Common preamble for all lab reports
\documentclass[a4paper, 12pt]{article}

% Language and encoding
\usepackage[T2A]{fontenc}
\usepackage[utf8]{inputenc}
\usepackage[ukrainian]{babel}

% Typography
\usepackage{microtype}
\usepackage{lmodern}

% Page geometry
\usepackage[margin=25mm, headheight=14pt]{geometry}

% Math
\usepackage{amsmath, amssymb}

% Graphics and figures
\usepackage{graphicx}
\usepackage{float}
\usepackage{subcaption}
\usepackage[justification=centering]{caption}

% Tables
\usepackage{tabularx}
\usepackage{booktabs}
\usepackage{array}

% Colors and hyperlinks
\usepackage{xcolor}
\usepackage[colorlinks=true, linkcolor=blue!60!black, urlcolor=blue!60!black, citecolor=blue!60!black]{hyperref}

% Lists
\usepackage{enumitem}
\setlist{nosep, leftmargin=*}

% Section formatting
\usepackage{titlesec}
\titleformat{\section}{\large\bfseries}{\thesection.}{0.5em}{}
\titleformat{\subsection}{\normalsize\bfseries}{\thesubsection.}{0.5em}{}

% Header/footer
\usepackage{fancyhdr}
\pagestyle{fancy}
\fancyhf{}
\fancyhead[L]{\small\leftmark}
\fancyhead[R]{\small\thepage}
\renewcommand{\headrulewidth}{0.4pt}
\fancypagestyle{plain}{\fancyhf{}\fancyfoot[C]{\thepage}\renewcommand{\headrulewidth}{0pt}}

% Placeholder image command
\newcommand{\placeholder}[2][0.8\textwidth]{%
  \begin{figure}[H]
    \centering
    \fbox{\parbox{#1}{\vspace{4cm}\centering\textit{#2}\vspace{4cm}}}
    \caption{#2}
  \end{figure}
}

% Placeholder for inline result values
\newcommand{\result}[1]{\textbf{\textcolor{red!70!black}{[#1]}}}

\usepackage{../common/labstyle}

\labnum{5}
\labtitle{Перемішування (Immersed Solid)}

\begin{document}
\maketitle
\tableofcontents
\newpage

% ===================================================================
\section{Вихідні дані}
% ===================================================================

\begin{table}[H]
  \centering
  \caption{Вихідні дані (варіант~13)}
  \begin{tabularx}{0.75\textwidth}{Xr}
    \toprule
    \textbf{Параметр} & \textbf{Значення} \\
    \midrule
    Тип мішалки                  & Якірна (Anchor) \\
    Швидкість обертання          & 400~об/хв ($\approx 41{,}9$~рад/с) \\
    Час дослідження              & 4~с \\
    Рідина                       & Water \\
    Метод моделювання            & Immersed Solid \\
    Тип розрахунку               & Transient \\
    \bottomrule
  \end{tabularx}
\end{table}

Геометрія резервуара --- внутрішній об'єм ємності з Практичної роботи~\textnumero 1.

\subsection{Параметри перемішування}

Часовий крок (36 кроків на оберт):
$\Delta t = 60/(\mathrm{RPM} \times 36) = 60/(400 \times 36) \approx 4{,}17 \times 10^{-3}$~с.

Число Рейнольдса: $\mathrm{Re_{mix}} = \rho N D^2 / \mu$.

Метод Immersed Solid --- тверде тіло ``занурюється'' в загальну сітку рідинного домену. Елементи, що перекриваються з тілом, отримують силові члени (source terms). Перевага: не потрібна окрема сітка для обертової зони.

% ===================================================================
\section{Послідовність виконання}
% ===================================================================

Робота виконана через ANSYS Workbench Journaling (\texttt{lab5.wbjn}) з солвером CFX.

\begin{lstlisting}[style=scheme]
1. CreateSystem("Fluid Flow (CFX)")
2. Import lab5.step (assembly: vessel_fluid + mixer,
   single file)
3. Enter Meshing
   |-- ID bodies by volume: vessel_fluid=largest,
   |                         mixer=smaller
   |-- Print body names (verify match CCL refs)
   |-- NS "Top"    -> vessel faces Z > 995mm (flat lid)
   |-- NS "Bottom" -> vessel faces Z < 0 (hemisphere)
   |-- NS "Wall"   -> remaining vessel faces (cylinder)
   |-- NS "Mixer"  -> all mixer body faces
   |-- Mesh 8mm -> mesh.png
4. Enter CFX-Pre
   |-- Import CCL:
   |   FluidDomain: Location=vessel_fluid, Water, k-eps
   |     Wall: No Slip | Top: Free Slip | Bottom: No Slip
   |   IMMERSED SOLID MixerSolid:
   |     Location=Mixer, Rotating 41.888 rad/s
   |     Axis = Coord 0.3 (Z-axis, geometry is Z-up)
   |   Transient: dt=4.17ms, T=4s (~960 steps)
   |   Output: P,V,TKE,TEF,WallShear every 24 steps
   |           -> 40 frames
   |   Parallel: -part 12
5. Solve
6. Enter CFX-Post
   |-- Plane: XZ at Y=0
   |-- Export: velocity, pressure, ssr, tke, tef,
   |           geometry
   '-- Export 40 frames -> ffmpeg -> animation.mp4
\end{lstlisting}

\subsection{Геометрія (CadQuery $\to$ STEP)}

Скрипт \texttt{geometry/lab5.py} будує два тіла: \textbf{vessel\_fluid} (внутрішній об'єм ємності) та \textbf{mixer} (якірна мішалка: U-труба + центральна вісь, радіус руків 220~мм, діаметр труби 20~мм). Експорт: \texttt{geometry/step/lab5.step} (assembly).

\placeholder{Геометрія резервуара з якірною мішалкою}

\placeholder{Розрахункова сітка (Mesh)}

% ===================================================================
\section{Результати розрахунку}
% ===================================================================

\subsection{Velocity}
\placeholder{Contour: розподіл швидкості (Velocity)}

\subsection{Pressure}
\placeholder{Contour: розподіл тиску (Pressure)}

\subsection{Shear Strain Rate}
\placeholder{Contour: Shear Strain Rate}

\subsection{Turbulence Kinetic Energy}
\placeholder{Contour: Turbulence Kinetic Energy}

\subsection{Turbulence Eddy Frequency}
\placeholder{Contour: Turbulence Eddy Frequency}

\subsection{Анімація}

Створено відеофайл \texttt{animation.mp4}, що демонструє процес перемішування рідини (4~с симуляції).

% ===================================================================
\section{Аналіз та висновок}
% ===================================================================

\result{Опис:
\begin{itemize}
  \item Структура течії: циркуляційні потоки, осьовий, радіальний та тангенціальний компоненти;
  \item Зони інтенсивного перемішування та зони застою;
  \item Розподіл TKE --- зони генерації турбулентності;
  \item Розподіл Shear Strain Rate --- зони максимального зсуву;
  \item Оцінка ефективності якірної мішалки при 400~об/хв;
  \item Загальний висновок.
\end{itemize}}

\end{document}
